% Catalogue of new phenomena and testable predictions
%
% This file is included from Part~III (\section{Predictions}).
% It summarises the genuinely new physical phenomena implied by the axiom
% $\Phi^2 = \Phi + 1$, organised by domain. Each item points back to the
% derivation in Part~I and the comparison with observation in Part~II.

\subsubsection*{Quantum sector}

\begin{itemize}[nosep,leftmargin=*]
\item \textbf{Universal coherence quality factor $Q$.}
  The theory predicts a universal ratio of coherent oscillations to decoherence,
  \[
    Q \;=\; \frac{\pi}{2\ln\varphi} \;\approx\; 3.264,
  \]
  derived from the breathing propagator and the axiom (Section~\ref{sec:quantum-predictions}).
  This sets an upper bound on the number of coherent cycles before damping at any scale.

\item \textbf{Coherence length $l_{\mathrm{coh}} \sim l_P \varphi^{Q}$.}
  Combining $Q$ with the Planck length yields a specific coherence length
  at which breathing structure becomes resolvable,
  \[
    l_{\mathrm{coh}} \;\sim\; l_P \varphi^{Q},
  \]
  providing a concrete quantum--gravity transition scale rather than a generic
  “Planck-scale neighbourhood” (Section~\ref{sec:quantum-predictions}).

\item \textbf{Casimir $\varphi$-excess and GW $\varphi$-ratio echoes.}
  The vacuum structure implies a Casimir force larger than the standard QFT
  prediction by a fixed factor of order~$\varphi$ and predicts $\varphi$-ratio
  echoes in gravitational-wave signals at specific frequency ratios
  (Section~\ref{sec:quantum-predictions} and the falsifiability discussion
  in Section~\ref{sec:verification-falsifiability}). Both are absent in the
  standard formalism and provide direct tests of the breathing vacuum.

\item \textbf{Breathing interior and absence of classical reality.}
  Signals whose difference/ratio transforms fail to converge to a stable
  recurrence are assigned to the breathing interior ($s \in (0,1)$).
  For such quantities the theory predicts an intrinsic absence of classical
  definiteness rather than a failure of numerical methods; the “lack of
  pattern” is itself a physical prediction (Section~\ref{sec:qm} and
  the breathing-summary discussion).
\end{itemize}

\subsubsection*{Cosmology and dark sector}

\begin{itemize}[nosep,leftmargin=*]
\item \textbf{Kiss redshift and onset of acceleration.}
  The twin-universe geometry fixes the deceleration--acceleration transition at
  the Kiss node,
  \[
    z_{\mathrm{kiss}} \;=\; \frac{1}{\varphi} \;\approx\; 0.618,
  \]
  with $q(z_{\mathrm{kiss}})=0$ and $\Omega_\Lambda/\Omega_m = \varphi^3/2$
  (Sections~\ref{sec:twin-lifecycle} and~\ref{sec:kiss-cosmology}).

\item \textbf{Vacuum energy $\rho_\Lambda/\rho_P = e^{-281}$.}
  At $k=4$ the effective number of modes $N_{\mathrm{eff}}=281$ is fixed by
  three independent spiral constructions; the resulting vacuum energy is
  (Section~\ref{sec:vacuum-energy})
  \[
    \frac{\rho_\Lambda}{\rho_P} \;=\; e^{-281} \;\approx\; 10^{-122.0},
  \]
  within one decade on a 123-decade span of the observed value.

\item \textbf{Dark-matter to baryon ratio $\Omega_{\mathrm{DM}}/\Omega_b
  = \varphi^2/\ln\varphi$.}
  The spectral decomposition of the partition function for the two branches
  yields
  \[
    \frac{\Omega_{\mathrm{DM}}}{\Omega_b}
      \;=\; \frac{\varphi^2}{\ln\varphi},
  \]
  matching the Planck value within $1.3\sigma$ without a free density parameter
  (Section~\ref{sec:dm-ratio}).

\item \textbf{Dark matter as geometry, not a particle.}
  The “dark matter” contribution is the coupling imprint of the $\psi$-mode
  from the Bed universe, required by $L_k = \varphi^k + \psi^k$.
  No new particle species (WIMP, axion) are introduced; direct and collider
  searches are predicted to remain null (Section~\ref{sec:cosmology-compatibility}).

\item \textbf{X17 boson and nuclear/gauge screening.}
  A protophobic boson with bare mass
  \[
    m_X^{\mathrm{bare}} = m_p \varphi^{-8}\bigl(1 - 2/L_6\bigr)
  \]
  and effective mass
  \[
    m_X^{\mathrm{eff}} = m_X^{\mathrm{bare}}\,\frac{L_8}{L_4^2}
  \]
  is predicted at $\sim 17~\mathrm{MeV}$, with coupling bounded by
  $(C(5)/C(4))^2 = (47/123)^2$.  The screening factor $L_8/L_4^2 = 47/49$
  is fixed by the depth chain and gauge structure (Section~\ref{sec:x17} and
  companion Paper~1).

\item \textbf{Dark fermions at $k=5$ and suppressed coupling.}
  The spiral structure at $k=4$ and $k=5$ partitions the $d=56$ matter
  dimension into three visible generations and two dark fermions ($F_5=5$),
  coupled to the Standard Model with strength
  \[
    \frac{C(5)}{C(4)} \;=\; \frac{47}{123}.
  \]
  This fixes the minimal dark sector content and its coupling hierarchy
  (Section~\ref{sec:dark-sector}).
\end{itemize}

\subsubsection*{Particle physics and unification}

\begin{itemize}[nosep,leftmargin=*]
\item \textbf{Proton decay lifetime.}
  The spiral GUT scale $M_{\mathrm{GUT}} = M_P\varphi^{-7}$, combined with the
  standard GUT lifetime scaling, yields
  \[
    \tau_p \;\sim\; 10^{36.5}\,\mathrm{yr},
  \]
  above current bounds and within the planned reach of Hyper-Kamiokande
  (Section~\ref{sec:proton-decay}).

\item \textbf{Neutrino masses from geometric suppression.}
  A seesaw-like suppression
  \[
    \frac{m_\nu}{m_D} \;=\; C(k) \;=\; \frac{2}{L_{2k}}
  \]
  arises from the branch-symmetric coupling, yielding an effective seesaw
  without an explicit heavy right-handed neutrino $M_R$
  (Section~\ref{sec:neutrino-mass}).
\end{itemize}

\subsubsection*{Black holes and gravity}

\begin{itemize}[nosep,leftmargin=*]
\item \textbf{Black holes as maximal-coupling contact points.}
  Black holes are identified with discrete contact points where
  $C(k)=2/L_{2k}$ reaches unity and the two vortex sheets come into full
  contact.  They are structural features of the double-vortex geometry,
  not products of gravitational collapse in a fixed background
  (Section~\ref{sec:bh-geometry}).

\item \textbf{PBH evaporation halts at $N = \dim = 133$.}
  The Hawking cascade terminates when the spectrum reaches $N = \dim = 133$,
  leaving a stable remnant of mass $M \approx 5.2\times 10^{-6}M_\odot$,
  rather than complete evaporation with a final burst
  (Section~\ref{sec:bh-mass}).

\item \textbf{Discrete black-hole mass spectrum and log$_\varphi$ peaks.}
  The mass spectrum
  \[
    M(N) \;=\; \frac{M_H(z_{\mathrm{kiss}})}{\varphi^N}
  \]
  predicts discrete attractor masses with deserts between them; rebinned
  black-hole mass catalogues in $\log_\varphi$ spacing should show clustering
  at the predicted $N$-levels (Section~\ref{sec:bh-predictions} and
  the black-hole prediction catalog).
\end{itemize}

